\section*{Chapter \ref{ch:b3:instrumental-variables} --- Instrumental Variables}

\begin{solution}{exo:b3:instrumental-variables:1}
\begin{enumerate}
\item The slope is an elasticity: a country whose predicted share is 20~\%
higher trades, on average, $1.2^{0.235} - 1 \approx 4.4$~\% more (the linear
approximation $0.235 \times 20 = 4.7$~\% is close).
\item The line predicts $50.0$~\% for a built share of $2.33$~\% (the left end
of the line) and about $60.2$~\% at $5.11$~\%. % figdata: first_stage_line
In logarithms Fiji lies $\ln(110.7/50.0) \approx 0.79$ above the line and the
United States $\ln(26.3/60.2) \approx -0.83$ below it: the United States is
slightly further away, on the other side. Large economies trade less with
others than geography predicts, because so much of their trade is internal.
\item No. The \omterm{def:b3:instrumental-variables:tsls}{first stage} measures relevance, the correlation of the instrument
with trade. The \omterm{def:b3:instrumental-variables:instrument}{exclusion restriction} is about the unobserved causes of output
per worker, which the chart does not show; a strong \omterm{def:b3:instrumental-variables:tsls}{first stage} and a direct
effect of geography on output can coexist.
\end{enumerate}
\end{solution}

\begin{solution}{exo:b3:instrumental-variables:2}
\begin{enumerate}
\item $\hat\beta^{\mathrm{W}} = (-0.00898)/(-0.1256) = 0.0715$. % ledger: b3iv:ak91.panelA
\item It is below OLS (0.0801) in the 1970 sample, whereas in the 1980 sample
IV (0.1020) was above OLS (0.0709). The two cohorts disagree on the sign of
the gap, which is a reminder that a gap of this size between IV and OLS is
within the IV estimates' sampling error (standard errors 0.0219 and 0.0239).
\end{enumerate}
\end{solution}

\begin{solution}{exo:b3:instrumental-variables:3}
\begin{enumerate}
\item By the omitted-variable formula, $\plim \hat\beta_1 = 0.06 + 0.12 \times
0.25 = 0.09$.
\item By $(0.09 - 0.06)/0.06 = 50$~\%.
\item $0.06$: the omitted variable biases OLS only if it enters the wage
equation ($\gamma \neq 0$) and is correlated with schooling.
\end{enumerate}
\end{solution}

\begin{solution}{exo:b3:instrumental-variables:4}
\begin{enumerate}
\item Relevance: people who grew up closer to a college get, on average, more
schooling. Exclusion: distance to a college is unrelated to the other causes
of adult wages and affects wages only through schooling.
\item Colleges are built in towns, so a short distance means an urban
childhood, with its better schools and denser job networks; and families who
value education may choose to live near a college, so distance is correlated
with parental background.
\item For the urban channel: compare the wages of people who grew up near and
far from a college among those who never attended any college. If distance
affected wages only through schooling, the gap within this group, conditional
on schooling, would be zero; a gap shows the other channel. For the family
channel: check whether distance predicts parents' education, which schooling
cannot have caused.
\end{enumerate}
\end{solution}

\begin{solution}{exo:b3:instrumental-variables:5}
\begin{enumerate}
\item \omterm{def:b3:instrumental-variables:tsls}{First stage} $0.40/0.25 = 1.6$; \omterm{def:b3:instrumental-variables:tsls}{reduced form} $0.06/0.25 = 0.24$; IV
estimate $0.06/0.40 = 0.15$, which is also $0.24/1.6$.
\item $\sqrt{0.36 \times 0.25/(1000 \times 0.16)} = \sqrt{0.0005625} = 0.0237$;
the interval is $0.15 \pm 1.96 \times 0.0237 = 0.15 \pm 0.0465$, from 0.104 to
0.196.
\item The fitted value is $\hat x_i = \hat a + \hat\pi z_i$ with $\hat\pi =
\widehat{\Cov}(z, x)/\widehat{\Var}(z)$. For a regression on a constant and one
variable, the $(2,2)$ entry of $(\hat X'\hat X)^{-1}$ is $1/\sum_i (\hat x_i -
\bar{\hat x})^2 = 1/(n \hat\pi^2 \widehat{\Var}(z))$ (variances with divisor
$n$). Substituting $\hat\pi$ gives $\widehat{\Var}(z)/(n\,\widehat{\Cov}(z,
x)^2)$, and multiplying by $\hat\sigma^2$ gives the formula.
\end{enumerate}
\end{solution}

\begin{solution}{exo:b3:instrumental-variables:6}
\begin{enumerate}
\item With classical error the OLS slope converges to $\lambda\beta$, where
$\lambda$ is the share of the variance of recorded schooling that is true
variance: $\lambda = 0.9$, so OLS converges to $0.9\beta$. A valid instrument
is uncorrelated with the measurement error, so IV converges to $\beta$.
\item OLS converges to $0.9\beta + 0.01$, which is below $\beta$ when $0.1\beta
> 0.01$, that is $\beta > 0.1$.
\item With unequal returns, IV estimates the average return among \omterm{def:b3:instrumental-variables:compliance}{compliers},
which can be larger than the average return that OLS (if it were unbiased)
would weight differently. Men kept in school by the leaving age may gain more
from a year than the average man.
\end{enumerate}
\end{solution}

\begin{solution}{exo:b3:instrumental-variables:7}
\begin{enumerate}
\item \omterm{def:b3:instrumental-variables:tsls}{First stage} $0.62 - 0.14 = 0.48$; \omterm{def:b3:instrumental-variables:tsls}{reduced form} $0.53 - 0.47 = 0.06$;
Wald estimate $0.06/0.48 = 0.125$. Among pupils who take tutoring if and only
if offered it, tutoring raises the chance of passing by 12.5 percentage points.
\item \omterm{def:b3:instrumental-variables:compliance}{Always-takers} $0.14$, \omterm{def:b3:instrumental-variables:compliance}{never-takers} $1 - 0.62 = 0.38$, \omterm{def:b3:instrumental-variables:compliance}{compliers} $0.48$.
\item \omterm{def:b3:instrumental-variables:compliance}{Monotonicity}: no pupil is tutored if not offered but untutored if
offered. A violation: a family that would have paid a private tutor, but
accepts the charity's offer, then finds the charity's sessions at an
impossible time and stops.
\end{enumerate}
\end{solution}

\begin{solution}{exo:b3:instrumental-variables:8}
The tested solution is \code{ex08_no_population.py}, in the chapter's
\code{exercises/} folder.
\begin{enumerate}
\item First-stage $F$ 16.7; OLS 0.833 (0.124); 2SLS 0.376 (0.451).
% output: ex08_no_population
\item The correlation of log population with the log built share is 0.092 and
with the log actual share $-0.544$. Population explains much of the actual
trade share and almost none of the built share. Controlling for it removes a
large part of the \omterm{def:b3:instrumental-variables:tsls}{first stage}'s residual variance without changing much the
variation of the instrument, so the $t$ statistic of the instrument, and its
$F$, rise.
\item Less plausible. Large countries trade less and may have different
productivity for reasons of size (larger markets, more varied industries);
without the control, those channels sit in the error. Since the built share is
nearly uncorrelated with population the damage is small here, but the control
is the safer specification.
\end{enumerate}
\end{solution}

\begin{solution}{exo:b3:instrumental-variables:9}
\begin{enumerate}
\item Conflict in year $t - 1$: rainfall in year $t$ cannot have caused it
through growth in year $t$. If rainfall in $t$ predicts conflict in $t - 1$,
rainfall is correlated with persistent causes of conflict.
\item The \omterm{def:b3:instrumental-variables:tsls}{reduced form} is the \omterm{def:b3:instrumental-variables:tsls}{first stage} times the effect of growth on
conflict, so with a weaker \omterm{def:b3:instrumental-variables:tsls}{first stage} the \omterm{def:b3:instrumental-variables:tsls}{reduced form} should be smaller
downstream of dams. Sarsons found that riots responded as much there as
elsewhere: rainfall reaches conflict through another channel, and the
\omterm{def:b3:instrumental-variables:instrument}{exclusion restriction} fails.
\item The \omterm{def:b3:instrumental-variables:tsls}{reduced form} is $\pi\beta + c$, so the IV estimate converges to
$(\pi\beta + c)/\pi = \beta + c/\pi$. The bias $c/\pi$ grows without bound as
the \omterm{def:b3:instrumental-variables:tsls}{first stage} weakens: the weaker the rainfall--growth link becomes, the more
a small direct effect of rain distorts the estimate.
\end{enumerate}
\end{solution}

\begin{solution}{exo:b3:instrumental-variables:10}
The tested solution is \code{ex10_weak_f5.py}, in the chapter's
\code{exercises/} folder.
\begin{enumerate}
\item Mean first-stage $F$ 5.06; median 2SLS 0.028; mean OLS 0.784; the 5~\%
test rejects the true value in 9.0~\% of the samples. % output: ex10_weak_f5
\item 13.0~\% of the samples have $F > 10$. Their median 2SLS estimate is
0.357, against $-0.059$ in the others. % output: ex10_weak_f5
A large sample \omterm{def:b3:instrumental-variables:tsls}{first stage}, $\widehat{\Cov}(z, x)$, is often large because
$\widehat{\Cov}(z, v)$ happens to be large; since $u$ and $v$ are correlated,
$\widehat{\Cov}(z, u)$ is then large too, with the sign of the OLS bias.
Selecting on the \omterm{def:b3:instrumental-variables:tsls}{first stage} therefore selects samples whose estimates lean
toward OLS, and a literature filtered on $F > 10$ is biased toward OLS.
\item An $F$ of 5 is below every critical value of the table, including the
5.53 needed for a size of at most 25~\%: the table cannot promise even that.
The 9.0~\% found here is for a correlation of 0.8 between $u$ and $v$; the
critical values guard against the worst correlation.
\end{enumerate}
\end{solution}

\begin{solution}{exo:b3:instrumental-variables:11}
\begin{enumerate}
\item By part (1) of \cref{prop:b3:instrumental-variables:tsls}, 2SLS is OLS of
$y$ on $[W, \hat x]$, where $\hat x = Wa + \hat\pi z$ is the first-stage fitted
value and $\hat\pi$ its coefficient on $z$. Let $M_W$ be the matrix that
residualises on $W$ and write $\tilde w = M_W w$. By the Frisch--Waugh--Lovell
theorem, the coefficient on $\hat x$ is $\tilde{\hat x}'\tilde y/\tilde{\hat
x}'\tilde{\hat x}$. Since $M_W W = 0$, $\tilde{\hat x} = \hat\pi\tilde z$, so
the coefficient is
\[
  \frac{\hat\pi\,\tilde z'\tilde y}{\hat\pi^2\,\tilde z'\tilde z} = \frac{1}{\hat\pi}\cdot\frac{\tilde z'\tilde y}{\tilde z'\tilde z}.
\]
By the same theorem, $\tilde z'\tilde y/\tilde z'\tilde z$ is the reduced-form
coefficient on $z$. The 2SLS coefficient is the reduced-form coefficient
divided by the first-stage coefficient.
\item If the reduced-form coefficient is zero, the ratio is zero, whatever the
value of $\hat\pi \neq 0$. An instrument that does not move the outcome cannot
produce an effect.
\end{enumerate}
\end{solution}

\begin{solution}{exo:b3:instrumental-variables:12}
\begin{enumerate}
\item As in the proof of \cref{thm:b3:instrumental-variables:late}, the
numerator is $\E[(D_i(1) - D_i(0))(Y_i(1) - Y_i(0))]$. The factor is 1 for
\omterm{def:b3:instrumental-variables:compliance}{compliers} and $-1$ for \omterm{def:b3:instrumental-variables:compliance}{defiers}, so the numerator is $\pi_c\Delta_c -
\pi_d\Delta_d$. The denominator is $\E[D_i(1) - D_i(0)] = \pi_c - \pi_d$.
\item $(0.3 \times 0.1 - 0.2 \times 0.5)/(0.3 - 0.2) = (0.03 - 0.10)/0.1 =
-0.7$. Every effect is positive and the ratio is negative, outside the range of
every individual effect: without \omterm{def:b3:instrumental-variables:compliance}{monotonicity}, the Wald ratio is not an
average effect of anyone.
\item Defendants assigned at random to judges: a judge who is lenient on
average may still be harsher than another judge towards some kinds of case, so
some defendants would be jailed by the lenient judge and freed by the strict
one.
\end{enumerate}
\end{solution}

\begin{solution}{pb:b3:instrumental-variables:lottery}
\textbf{1.} $1{,}800/4{,}000 = 0.45$ enrolled. The enrolled mean is
$(1{,}400 \times 10.2214 + 400 \times 10.40)/1{,}800 = 18{,}469.96/1{,}800 =
10.261$; the non-enrolled mean is $(600 \times 9.80 + 1{,}600 \times
9.89375)/2{,}200 = 21{,}710/2{,}200 = 9.868$.

\textbf{2.} $10.261 - 9.868 = 0.393$.

\textbf{3.} Applicants who enrol differ from those who do not in ways that
matter for earnings: those who enrol after losing (\omterm{def:b3:instrumental-variables:compliance}{always-takers}) are
determined or well-connected, and those who decline a place they won
(\omterm{def:b3:instrumental-variables:compliance}{never-takers}) may have good alternatives or poor prospects. And the gain from
college may differ between those who choose it and those who do not.

\textbf{4.} No. Grades capture part of what the applicants differ in, but the
selection runs on motivation, family resources and alternatives, which grades
do not measure; the residual difference is still not the effect of college.

\textbf{5.} $\eu^{0.393} - 1 = 0.48$: the prospectus claims earnings 48~\%
higher.

\textbf{6.} Relevance: winning raises the chance of enrolling. Exclusion:
winning affects earnings only through enrolment, not, for instance, through
the encouragement of having won. The lottery assigns winning by chance, so
winners and losers have the same distribution of potential outcomes and
potential enrolments: independence holds by design.

\textbf{7.} \omterm{def:b3:instrumental-variables:tsls}{First stage} $1{,}400/2{,}000 - 400/2{,}000 = 0.70 - 0.20 = 0.50$;
\omterm{def:b3:instrumental-variables:tsls}{reduced form} $10.095 - 9.995 = 0.100$.

\textbf{8.} $0.100/0.50 = 0.200$.

\textbf{9.} $\sqrt{0.000105 + 0.00008} = \sqrt{0.000185} = 0.0136$; $t =
0.50/0.0136 = 36.8$ and $F = t^2 \approx 1{,}350$. Not weak: far above every
threshold.

\textbf{10.} $0.60 \sqrt{1/2{,}000 + 1/2{,}000} = 0.60 \times 0.0316 = 0.0190$.

\textbf{11.} $0.0190/0.50 = 0.0380$.

\textbf{12.} $0.200 \pm 1.96 \times 0.0380 = 0.200 \pm 0.074$, from 0.126 to
0.274.

\textbf{13.} \omterm{def:b3:instrumental-variables:compliance}{Always-takers}: the enrolled share among losers, 0.20.
\omterm{def:b3:instrumental-variables:compliance}{Never-takers}: the non-enrolled share among winners, 0.30. \omterm{def:b3:instrumental-variables:compliance}{Compliers}: the \omterm{def:b3:instrumental-variables:tsls}{first
stage}, 0.50.

\textbf{14.} The enrolled losers are all \omterm{def:b3:instrumental-variables:compliance}{always-takers}: their mean of $Y(1)$ is
10.40. The non-enrolled winners are all \omterm{def:b3:instrumental-variables:compliance}{never-takers}: their mean of $Y(0)$ is
9.80. (The \omterm{def:b3:instrumental-variables:compliance}{always-takers}' $Y(0)$ and the \omterm{def:b3:instrumental-variables:compliance}{never-takers}' $Y(1)$ are never seen.)

\textbf{15.} Among winners, $0.7 \times 10.2214$ is the share-weighted mean
of $Y(1)$ of the \omterm{def:b3:instrumental-variables:compliance}{always-takers} (share 0.2, mean 10.40) and the \omterm{def:b3:instrumental-variables:compliance}{compliers}
(share 0.5), by independence. So the \omterm{def:b3:instrumental-variables:compliance}{compliers}' mean is $(0.7 \times 10.2214 -
0.2 \times 10.40)/0.5 = (7.155 - 2.080)/0.5 = 10.150$.

\textbf{16.} Among losers, the non-enrolled are \omterm{def:b3:instrumental-variables:compliance}{never-takers} (share 0.3, mean
9.80) and \omterm{def:b3:instrumental-variables:compliance}{compliers} (share 0.5): $(0.8 \times 9.89375 - 0.3 \times 9.80)/0.5
= (7.915 - 2.940)/0.5 = 9.950$.

\textbf{17.} $10.150 - 9.950 = 0.200$, the Wald estimate.

\textbf{18.} Without college the \omterm{def:b3:instrumental-variables:compliance}{compliers} earn more than the \omterm{def:b3:instrumental-variables:compliance}{never-takers}
(9.95 against 9.80), and with college less than the \omterm{def:b3:instrumental-variables:compliance}{always-takers} (10.15
against 10.40). Enrolment rises with earnings potential: the people who enrol
whatever happens are the highest earners, those who never enrol the lowest.
The naive comparison of Part I adds that ordering to the effect of college,
which is why it is about twice the Wald estimate.

\textbf{19.} \omterm{def:b3:instrumental-variables:tsls}{First stage} $0.26 - 0.20 = 0.06$; standard error $\sqrt{0.1924/600
+ 0.16/600} = \sqrt{0.000587} = 0.0242$; $t = 0.06/0.0242 = 2.48$ and $F = 2.48^2
\approx 6.2$.

\textbf{20.} $0.02/0.06 = 0.333$.

\textbf{21.} $0.60\sqrt{2/600} = 0.60 \times 0.0577 = 0.0346$; the Wald
estimate's standard error is about $0.0346/0.06 = 0.577$.

\textbf{22.} The difference is $0.333 - 0.200 = 0.133$, with standard error
$\sqrt{0.0380^2 + 0.577^2} = 0.578$; the $z$ statistic is 0.23. The two
estimates are not distinguishable: the second campus's is too imprecise to
agree or disagree with anything.

\textbf{23.} Weak: an $F$ of about 6.2 is below the rule of thumb of 10 and
below the Stock--Yogo critical values for a size of 10~\%, 15~\% and 20~\%;
its $t$-test could be badly oversized, and the normal approximation behind
Question~21 is itself unreliable. Pooling adds almost no information, because
2SLS weights each instrument by its first-stage strength, and adding a \omterm{def:b3:instrumental-variables:weak}{weak
instrument} moves 2SLS toward OLS in finite samples. Report the main campus,
and the second campus separately as a weak, uninformative estimate.

\textbf{24.} The two lotteries may move different people: the \omterm{def:b3:instrumental-variables:compliance}{compliers} at a
small campus, those who enrol only if they win a place there, may have
different gains from those at the main campus. Two valid LATEs can differ
without either being wrong. The named result: for applicants to the main
campus who enrol if and only if they win a place, college raises log earnings
at thirty by 0.200, about 22~\% ($\eu^{0.2} - 1$), with a 95~\% interval of
0.126 to 0.274 log points.
\end{solution}
